\section{Conclusion}
\label{sec:conclusion}

We tested whether a substantive belief-based persona -- one that names epistemic commitments rather than demographic identity or Big-Five traits -- produces more collaborative friction in production LLMs than the standard baselines (no persona, ``helpful assistant,'' role-based, constitutional). On the primary probe of productive correction of confidently-wrong user claims, the belief persona reduced sycophantic agreement by \(21.3\) percentage points on \gptone (\(p=0.0018\)) and \(18.8\) percentage points on \gptmini (\(p=0.007\)), with both effects surviving Holm correction. It was also the only persona whose surface style stuck to a held-out reference (Cohen's \(d \geq 1.27\)). Demographic and shallow personas produced no significant behavioral effect, replicating \citet{han2025personality} for those persona families. The Personality Illusion is therefore a narrower phenomenon than originally framed: persona injection \emph{is} action-controlling when the persona names what counts as evidence and what to do under pressure.

\para{Key takeaway.} Distinct identity, in the substantive sense of named epistemic commitments, is a necessary precondition for principled friction in current LLMs. The popular ``you are a helpful assistant'' framing is not neutral; it is yielding.

\para{Future work.} Six follow-ups directly extend this study. (i) Multi-turn extension à la \citet{liu2025truthdecay}: does the belief persona hold up after two or three rounds of user pushback? (ii) Cross-vendor replication on Claude, Gemini, and open-weight Llama-3 to test whether the effect is specific to OpenAI fine-tuning. (iii) Persona-text sensitivity: three variants each of the belief and constitutional persona, to separate category effect from wording effect. (iv) Goal-directed friction in \textsc{Sotopia}-style social scenarios, where the agent has a goal it must press for rather than fold for harmony. (v) Training-time persona induction via small DPO runs on belief-vs-no-persona pairs from \etwo, testing whether the prompt-only effect transfers to weights. (vi) Persona-induced bias check: re-evaluate the belief persona on reasoning benchmarks (\eg MMLU) to confirm that its sycophancy reduction does not come at the cost of new biased reasoning, as documented for other persona families by \citet{gupta2023bias}.
